\paragraph{\trSuspension{}操作。}\label{sub:shift-operation}
对任意 $\sigma\in\ZZ$，\trLinear\trLagrange\trBrane{}上存在一个 $\sigma$
重的\textbf{\trSuspension{}操作}（$\sigma$-fold shift operation），它把
$\Lambda^\#=(\Lambda,\alpha^\#,P^\#)$ 映为
\begin{equation}\label{eq:shift-operation}
S^\sigma\Lambda^\#=(\Lambda,\alpha^\#-\sigma,
P^\#\otimes\lambda^{\mathrm{top}}(\Lambda)^{\otimes\sigma}).
\end{equation}
当 $\sigma=1$ 时就简写作 $S\Lambda^\#$。关键性质是：
\begin{lemma}\label{thm:shift-index-orientation}
对任意一组满足 $\Lambda_0\pitchfork\Lambda_1$ 的\trLinear\trLagrange\trBrane{}
$\Lambda_0^\#,\Lambda_1^\#$，
\begin{equation}\label{eq:shift-index-orientation}
\begin{aligned}
i(\Lambda_0^\#,S^\sigma\Lambda_1^\#)&=i(\Lambda_0^\#,\Lambda_1^\#)-\sigma,\\
o(\Lambda_0^\#,S^\sigma\Lambda_1^\#)&\cong o(\Lambda_0^\#,\Lambda_1^\#).
\end{aligned}
\end{equation}
\end{lemma}
\begin{proof}
关于\trIndex{}的陈述是直截了当的，故我们只考虑\trOrientation{}空间的部分。为简单计，
设 $\sigma=1$。在\trLagrangianGrassmannian{}中，取一条从 $\Lambda_0$ 到 $\Lambda_1$ 的路径
$\lambda$，再取一个从 $\Lambda_1$ 到其自身的\trLoop{} $\rho$，并把两者的连接记作
$\lambda\cdot\rho$。\trUpperHalfPlane{}与\trClosed\trDisk{}上各有关联的算子 $D_{H,\lambda}$、
$D_{H,\lambda\cdot\rho}$ 与 $D_{D,\rho}$，它们之间的关系如下：若作一个%
\trBoundaryConnectedSum，把某一点 $s\in\partial H$（$s\gg 0$）与 $1\in\partial D$
等同起来，则据此把 $D_{H,\lambda}$ 与 $D_{D,\rho}$ \trGluing{}起来而得的算子可以等同于
$D_{H,\lambda\cdot\rho}$。
现在设 $\lambda$ 具有如下性质：函数 $\alpha_V(\lambda(t))$ 有一个实值\trLift，其在
两端点处的值为 $\alpha_0^\#$、$\alpha_1^\#$。再设 $\mu(\rho)=-1$，于是若把我们的实值%
\trLift{}扩张为 $\lambda\cdot\rho$ 的一个实值\trLift，则在新端点处的值将是
$\alpha_1^\#-1$。其次，考虑关联丛 $F_\lambda\to[0;1]$ 与 $\widetilde{F}_\rho\to[0;1]$
（后一记号的选取是为了与\cref{sub:pin-structures}中的记号保持一致）。我们可以为它们
配备 \trPinStructure{} $P_\lambda^\#$、$\widetilde{P}_\rho^\#$ 以及\trIsomorphism
\begin{equation}\label{eq:shift-proof-pin-isos}
\left\{
\begin{aligned}
(P_\lambda^\#)_0&\cong P_0^\#,\\
(P_\lambda^\#)_1&\cong P_1^\#,
\end{aligned}
\right.
\qquad
\left\{
\begin{aligned}
(\widetilde{P}_\rho^\#)_0&\cong P_1^\#,\\
(\widetilde{P}_\rho^\#)_1&\cong P_1^\#\otimes\lambda^{\mathrm{top}}(\Lambda_1).
\end{aligned}
\right.
\end{equation}
把这两者组合起来——利用\eqref{eq:shift-proof-pin-isos} 中的第二与第三个\trIsomorphism——
便得到 $F_{\lambda\cdot\rho}$ 上的一个 \trPinStructure，它在两端点处分别带有到
$P_0^\#$ 与 $P_1^\#\otimes\lambda^{\mathrm{top}}(\Lambda_1)$ 的\trCanonical\trIsomorphism。
由于\cref{thm:halfplane-index-determinant}以及\trOrientation{}空间的定义，$F_\lambda$
与 $F_{\lambda\cdot\rho}$ 上 \trPinStructure{}的选取确定了\trIsomorphism
$\det(D_{H,\lambda})\cong o(\Lambda_0^\#,\Lambda_1^\#)$ 与
$\det(D_{H,\lambda\cdot\rho})\cong o(\Lambda_0^\#,S\Lambda_1^\#)$。另一方面，
$\widetilde{P}_\rho^\#$ 是\eqref{eq:twisted-pin-structure} 意义下的一个\trTwisted{}
\trPinStructure，由\cref{thm:twisted-pin-determinant-iso}，它给出一个\trIsomorphism
$\det(D_{D,\rho})\cong\lambda^{\mathrm{top}}\Lambda_1$。鉴于\eqref{eq:boundary-gluing}，
我们于是有
\begin{align*}
o(\Lambda_0^\#,S\Lambda_1^\#)\otimes\lambda^{\mathrm{top}}(\Lambda_1)
&\cong\det(D_{H,\lambda\cdot\rho})\otimes\lambda^{\mathrm{top}}(\Lambda_1)\\
&\cong\det(D_{H,\lambda})\otimes\det(D_{D,\rho})
\cong o(\Lambda_0^\#,\Lambda_1^\#)\otimes\lambda^{\mathrm{top}}(\Lambda_1).\qedhere
\end{align*}
\end{proof}
\begin{remark}\label{thm:shift-alternative-proof}
与其诉诸\cref{thm:twisted-pin-determinant-iso}，也可以照搬其证明的一部分。具体地说，
取 $\rho(s)=\Lambda_{1,+}\oplus\exp(-\pi s I_V)\Lambda_{1,-}$，其中
$\Lambda_1=\Lambda_{1,+}\oplus\Lambda_{1,-}$ 是某个满足 $\dim(\Lambda_{1,-})=1$ 的%
\trSplitting。看一看\eqref{eq:halfplane-index-formula} 便可发现，$(\delta_H)_{\lambda\cdot\rho}$
与 $(\delta_H)_\lambda$ 相差一个增补进来的张量因子 $\Lambda_{1,-}$。而\trTwisted{}
\trPinStructure{}的选取恰好对应于同一空间的一个\trOrientation，这就使人得以认定
$(\delta_H)_{\lambda\cdot\rho}\cong(\delta_H)_\lambda$。这一论证更为初等，因为它避开了%
\trIndex{}理论；然而原有的证明对某些目的来说更为方便，比如理解
\eqref{eq:shift-index-orientation} 与\eqref{eq:general-surface-index-formula} 中诸%
\trIsomorphism{}之间的相互作用。
\end{remark}
